% ============================================================
%  EXAMEN — APLICACIONES DE LA DERIVADA A LA ECONOMÍA
%  90 minutos · 20 puntos
%  Secciones: 1) Monotonía  2) Máximos/mínimos (2.ª derivada)
%             3) Concavidad e inflexión  4) Optimización aplicada
%  Autor: Ing. Gabriel Astudillo
%  Compilar:  latexmk -xelatex examen_aplicaciones_derivadas.tex
%  Nota: la Clave de Respuestas está envuelta en \ifmwclaves; la versión
%        _scr (sin clave) se genera con \mwclavesoff.
% ============================================================

\documentclass[12pt, letterpaper]{article}

% ── Paquetes ─────────────────────────────────────────────────
\usepackage{mathwizards-guia}
\mwguia{Examen — Aplicaciones de la Derivada}{Ing. Gabriel Astudillo $\cdot$ Matemática Aplicada a la Economía}

\begin{document}
% ============================================================

% ── Portada / Título ─────────────────────────────────────────
\mwportada{Examen de Aplicaciones de la Derivada}{Cálculo Diferencial aplicado a la Economía $\cdot$ 90 minutos}{Monotonía, extremos, concavidad y optimización}

\vspace{6pt}
\noindent\textbf{Nombre:} \rule{0.5\linewidth}{0.4pt} \quad \textbf{Fecha:} \rule{0.18\linewidth}{0.4pt}

\vspace{2pt}
\noindent\small Tiempo: 90 minutos \quad$\bullet$\quad Total: 20 puntos \quad$\bullet$\quad Muestra \emph{todos} los pasos y justifica con las derivadas.

\vspace{10pt}

% ============================================================
\section{Monotonía: crecimiento y decrecimiento (4 puntos)}
% ============================================================

\begin{enumerate}[leftmargin=*]
  \item \textbf{(2 pts)} Para $f(x) = x^{3} - 3x^{2} - 9x + 5$:
        \begin{enumerate}[label=\alph*)]
          \item Halla los intervalos de \textbf{crecimiento} y \textbf{decrecimiento}.
          \item Clasifica los \textbf{extremos relativos} (máximo y mínimo).
        \end{enumerate}

  \item \textbf{(2 pts)} Para $g(x) = \dfrac{x}{x^{2} + 1}$:
        \begin{enumerate}[label=\alph*)]
          \item Determina los intervalos de \textbf{monotonía}.
          \item Clasifica los \textbf{extremos relativos}.
        \end{enumerate}
\end{enumerate}

% ============================================================
\section{Máximos y mínimos con la segunda derivada (4 puntos)}
% ============================================================

\begin{enumerate}[leftmargin=*]
  \setcounter{enumi}{2}
  \item \textbf{(2 pts)} Para $f(x) = x^{3} - 6x^{2} + 9x + 1$, halla los puntos críticos y clasifícalos aplicando el \textbf{criterio de la segunda derivada}.

  \item \textbf{(2 pts)} Para $h(x) = 2x^{3} - 3x^{2} - 12x + 5$, halla los puntos críticos y clasifícalos aplicando el \textbf{criterio de la segunda derivada}.
\end{enumerate}

% ============================================================
\section{Concavidad y puntos de inflexión (4 puntos)}
% ============================================================

\begin{enumerate}[leftmargin=*]
  \setcounter{enumi}{4}
  \item \textbf{(2 pts)} Para $f(x) = x^{4} - 6x^{2} + 2x$, determina los intervalos de \textbf{concavidad} y los \textbf{puntos de inflexión}.

  \item \textbf{(2 pts)} Para $g(x) = 3x^{4} - 4x^{3}$, determina los intervalos de \textbf{concavidad} y los \textbf{puntos de inflexión}.
\end{enumerate}

% ============================================================
\section{Optimización aplicada a la economía (8 puntos)}
% ============================================================

\begin{enumerate}[leftmargin=*]
  \setcounter{enumi}{6}
  \item \textbf{(4 pts) Maximización del beneficio.}
        Una empresa produce y vende $q$ unidades de un artículo. El \textbf{precio} por unidad está dado por la demanda
        \[
          p(q) = 600 - 2q
        \]
        y el \textbf{costo total} de producción es
        \[
          C(q) = 3q^{2} + 100q + 2000.
        \]
        \begin{enumerate}[label=\alph*)]
          \item Escribe la función \textbf{ingreso} $I(q)$ y la función \textbf{beneficio} $B(q)$.
          \item Halla la cantidad $q$ que \textbf{maximiza el beneficio}. Justifica con la segunda derivada.
          \item Determina el \textbf{precio} $p$ y el \textbf{beneficio máximo}.
        \end{enumerate}

  \item \textbf{(4 pts) Minimización del costo medio.}
        El costo total de producir $q$ unidades de un bien es
        \[
          C(q) = 2q^{2} + 32q + 200.
        \]
        \begin{enumerate}[label=\alph*)]
          \item Escribe la función de \textbf{costo medio} $\overline{C}(q) = \dfrac{C(q)}{q}$.
          \item Halla la cantidad $q$ que \textbf{minimiza el costo medio}. Justifica con la segunda derivada.
          \item Determina el \textbf{costo medio mínimo}.
        \end{enumerate}
\end{enumerate}

\vspace{12pt}
\begin{center}
  {\Large\bfseries\color{mwprimario} ¡Mucho éxito!}\\[4pt]
  {\large\color{gray} Ordena tu procedimiento, revisa los signos y no olvides justificar cada conclusión.}
\end{center}

% ============================================================
\ifmwclaves
  \section{Clave de Respuestas}
  % ============================================================

  \subsection*{Sección 1 · Monotonía: crecimiento y decrecimiento}

  \begin{enumerate}[leftmargin=*, label=\textbf{\arabic*.}]
    \item $f(x) = x^{3} - 3x^{2} - 9x + 5$.
    \[
      f'(x) = 3x^{2} - 6x - 9 = 3\,(x - 3)(x + 1)
      \;\Longrightarrow\; x = -1,\; x = 3
    \]
    \textbf{Signo de $f'$:} $+$ en $(-\infty,-1)$, $-$ en $(-1,3)$, $+$ en $(3,\infty)$.
    \begin{itemize}
      \item \textbf{Crece} en $(-\infty,-1) \cup (3,\infty)$.
      \item \textbf{Decrece} en $(-1,3)$.
      \item \textbf{Máximo relativo} en $x = -1$: $f(-1) = 10$.
      \item \textbf{Mínimo relativo} en $x = 3$: $f(3) = -22$.
    \end{itemize}

    \item $g(x) = \dfrac{x}{x^{2} + 1}$. Regla del cociente:
    \[
      g'(x) = \frac{(1)(x^{2}+1) - x(2x)}{(x^{2}+1)^{2}}
             = \frac{1 - x^{2}}{(x^{2}+1)^{2}}
      \;\Longrightarrow\; x = \pm 1
    \]
    El denominador es siempre positivo, así que el signo lo da $1-x^{2}$: $-$ en $(-\infty,-1)$, $+$ en $(-1,1)$, $-$ en $(1,\infty)$.
    \begin{itemize}
      \item \textbf{Crece} en $(-1,1)$.
      \item \textbf{Decrece} en $(-\infty,-1) \cup (1,\infty)$.
      \item \textbf{Máximo relativo} en $x = 1$: $g(1) = \dfrac{1}{2}$.
      \item \textbf{Mínimo relativo} en $x = -1$: $g(-1) = -\dfrac{1}{2}$.
    \end{itemize}
  \end{enumerate}

  \subsection*{Sección 2 · Máximos y mínimos con la segunda derivada}

  \begin{enumerate}[leftmargin=*, label=\textbf{\arabic*.}]
    \item $f(x) = x^{3} - 6x^{2} + 9x + 1$.
    \[
      f'(x) = 3x^{2} - 12x + 9 = 3\,(x - 1)(x - 3)
      \qquad
      f''(x) = 6x - 12
    \]
    Puntos críticos: $x = 1$ y $x = 3$.
    \begin{itemize}
      \item $f''(1) = -6 < 0 \;\Rightarrow\;$ \textbf{máximo} en $x = 1$: $f(1) = 5$.
      \item $f''(3) = 6 > 0 \;\Rightarrow\;$ \textbf{mínimo} en $x = 3$: $f(3) = 1$.
    \end{itemize}

    \item $h(x) = 2x^{3} - 3x^{2} - 12x + 5$.
    \[
      h'(x) = 6x^{2} - 6x - 12 = 6\,(x - 2)(x + 1)
      \qquad
      h''(x) = 12x - 6
    \]
    Puntos críticos: $x = -1$ y $x = 2$.
    \begin{itemize}
      \item $h''(-1) = -18 < 0 \;\Rightarrow\;$ \textbf{máximo} en $x = -1$: $h(-1) = 12$.
      \item $h''(2) = 18 > 0 \;\Rightarrow\;$ \textbf{mínimo} en $x = 2$: $h(2) = -15$.
    \end{itemize}
  \end{enumerate}

  \subsection*{Sección 3 · Concavidad y puntos de inflexión}

  \begin{enumerate}[leftmargin=*, label=\textbf{\arabic*.}]
    \item $f(x) = x^{4} - 6x^{2} + 2x$.
    \[
      f'(x) = 4x^{3} - 12x + 2
      \qquad
      f''(x) = 12x^{2} - 12 = 12\,(x - 1)(x + 1)
      \;\Longrightarrow\; x = \pm 1
    \]
    \begin{itemize}
      \item \textbf{Cóncava hacia arriba} en $(-\infty,-1) \cup (1,\infty)$ (donde $f'' > 0$).
      \item \textbf{Cóncava hacia abajo} en $(-1,1)$ (donde $f'' < 0$).
    \end{itemize}
    Como $f''$ cambia de signo en $x=-1$ y en $x=1$:
    \[
      f(-1) = -7 \;\Rightarrow\; \text{inflexión en } (-1,-7);
      \qquad
      f(1) = -3 \;\Rightarrow\; \text{inflexión en } (1,-3).
    \]

    \item $g(x) = 3x^{4} - 4x^{3}$.
    \[
      g'(x) = 12x^{3} - 12x^{2}
      \qquad
      g''(x) = 36x^{2} - 24x = 12x\,(3x - 2)
      \;\Longrightarrow\; x = 0,\; x = \tfrac{2}{3}
    \]
    \begin{itemize}
      \item \textbf{Cóncava hacia arriba} en $(-\infty,0) \cup \left(\tfrac{2}{3},\infty\right)$ (donde $g'' > 0$).
      \item \textbf{Cóncava hacia abajo} en $\left(0,\tfrac{2}{3}\right)$ (donde $g'' < 0$).
    \end{itemize}
    Ambos puntos son de inflexión:
    \[
      g(0) = 0 \;\Rightarrow\; (0,0);
      \qquad
      g\!\left(\tfrac{2}{3}\right) = -\tfrac{16}{27} \;\Rightarrow\; \left(\tfrac{2}{3}, -\tfrac{16}{27}\right).
    \]
  \end{enumerate}

  \subsection*{Sección 4 · Optimización aplicada a la economía}

  \begin{enumerate}[leftmargin=*, label=\textbf{\arabic*.}]
    \item \textbf{Maximización del beneficio.} Con $p(q) = 600 - 2q$ y $C(q) = 3q^{2} + 100q + 2000$.

    \textbf{a)} Ingreso y beneficio:
    \[
      I(q) = p(q)\cdot q = (600 - 2q)\,q = 600q - 2q^{2}
    \]
    \[
      B(q) = I(q) - C(q)
           = \big(600q - 2q^{2}\big) - \big(3q^{2} + 100q + 2000\big)
           = -5q^{2} + 500q - 2000
    \]

    \textbf{b)} Punto crítico y justificación:
    \[
      B'(q) = -10q + 500 = 0 \;\Longrightarrow\; q = 50
      \qquad
      B''(q) = -10 < 0
    \]
    Como $B''(50) < 0$, en $q = 50$ hay un \textbf{máximo}: se deben producir y vender \textbf{50 unidades}.

    \textbf{c)} Precio y beneficio máximo:
    \[
      p(50) = 600 - 2\cdot 50 = 500
    \]
    \[
      I(50) = 600(50) - 2(50)^{2} = 25\,000
      \qquad
      C(50) = 3(50)^{2} + 100(50) + 2000 = 14\,500
    \]
    \[
      B_{\max} = B(50) = -5(2500) + 500(50) - 2000 = 10\,500
    \]
    \textbf{Respuesta:} $q = 50$ unidades, precio $p = 500$ y beneficio máximo $= 10\,500$.

    \item \textbf{Minimización del costo medio.} Con $C(q) = 2q^{2} + 32q + 200$.

    \textbf{a)} Costo medio:
    \[
      \overline{C}(q) = \frac{C(q)}{q} = \frac{2q^{2} + 32q + 200}{q} = 2q + 32 + \frac{200}{q}
    \]

    \textbf{b)} Punto crítico y justificación:
    \[
      \overline{C}'(q) = 2 - \frac{200}{q^{2}} = 0
      \;\Longrightarrow\; q^{2} = 100
      \;\Longrightarrow\; q = 10 \quad (q > 0)
    \]
    \[
      \overline{C}''(q) = \frac{400}{q^{3}} > 0 \quad (q > 0)
    \]
    Como $\overline{C}''(10) > 0$, en $q = 10$ hay un \textbf{mínimo}: conviene producir \textbf{10 unidades}.

    \textbf{c)} Costo medio mínimo:
    \[
      \overline{C}(10) = 2(10) + 32 + \frac{200}{10} = 20 + 32 + 20 = 72
    \]
    \textbf{Respuesta:} $q = 10$ unidades y costo medio mínimo de $72$ por unidad.

  \end{enumerate}

  \vspace{10pt}
  \begin{cuadrosolucion}
    \textbf{Baremo rápido.} Sección 1: 4 pts \quad$\bullet$\quad Sección 2: 4 pts \quad$\bullet$\quad Sección 3: 4 pts \quad$\bullet$\quad Sección 4: 8 pts. Total: \textbf{20 puntos}. Se otorga puntaje parcial por el procedimiento correcto aunque el resultado final tenga un error aritmético.
  \end{cuadrosolucion}
\fi

\end{document}
